Generated game events can satisfy a syntactic schema and cite known facts while requesting inaccessible objects, unauthorized state transitions, or premature secret disclosure. Grounding and plausible dialogue \cite{S03_weir2024ontologically,S20_ashby2023personalized} supply candidate content, yet the game engine must still decide which state transitions to authorize.

\system{} enforces this decision at an external proposal-to-commit boundary (Fig.~\ref{fig:architecture}). Models and adapters submit typed candidates to a deterministic gate that checks action policies, preconditions, object reachability, NPC knowledge, disclosure permissions, effects, and quest-stage monotonicity. When a check fails, the gate returns structured errors for bounded repair; repair modifies only the candidate rather than canonical state, and mutation occurs only after same-state revalidation. A read-only eight-role ontology view links source declarations to these checks and public decisions without write authority, while hash-linked transition records support semantic replay and assignment-complete accounting.

Conceptually, this contract goes beyond factual grounding by separating an NPC's \emph{epistemic knowledge boundary} (facts an actor knows) from its \emph{disclosure authorization policy} (facts an actor may reveal at the current quest stage). Methodologically, our evaluation separates a zero-LLM-call deterministic verification layer (gate checks, guided repair, and engine replay) from a bounded stochastic proposer layer with explicit failure retention.

Across deterministic audits, a schema-only baseline committed 260/416 reference-invalid proposals across 32 finite-graph configurations, whereas all semantic-gate arms recorded \textbf{0/416} unsafe commits, 0/156 false rejections, and 128/260 $\rho$ repairs, while 480 $R{\le}64, d{\le}5$ procedural Datalog checks (\textbf{P9}) established 100.0\% oracle parity at $\le$7.12\,$\mu$s median latency vs.\ 46.67\% flat and 60.00\% 2-hop ($Q{=}1100.19$). Extending Theorems~1--4 for invariants \textbf{I1--I4}, attested direct HTTPS Anthropic Claude~5.5 API evaluations across Haiku, Sonnet, and Opus~5.5 confirm \textbf{0} unsafe commits across 105 independent 5-turn closed-loop trajectories (\textbf{P15}: 15/15 zero-leak vs.\ Safe-RL~\cite{alshiekh2018shield} 1/15 and NeMo-style~\cite{rebedea2023nemo} 1/15 proxy rails, 15/15 hybrid safe goals), 60 prospective factorial episodes (\textbf{P16A}: \texttt{ontology\_multihop} 30/30 zero-leak and 22 repairs vs.\ 22/30 \texttt{general\_rule}), 21 Godot~4.7.1 closed-loop episodes (\textbf{P16B}: 21/21 exact replays, 6/6 \texttt{SealedLighthouseMachine} terminal hashes), and 540 unified EN/KO escape checks (\textbf{P17}: $\kappa{=}0.8813$, 0\% FRR), alongside cross-family CLI tracks (\textbf{P11}--\textbf{P14}: $(a,E)$ repair 27/27 vs.\ 26/27 at 0 extra tokens), \texttt{DATASHEET\_v5.md}~\cite{gebru2021datasheets}, Croissant~1.0 v5~\cite{akhtar2024croissant} (26/26 SHA-256 digests), and quarantined \texttt{TR01--TR06} ($N{=}0$, \texttt{AWAITING\_HUMAN\_RATERS}).

\section{Related Work and Research Gap}
Grounding, memory, planning, and role-based agents \cite{S03_weir2024ontologically,S20_ashby2023personalized,S07_park2023generative,S08_shao2023character}, together with player and dialogue systems \cite{S21_akoury2023player,S22_hochreiter2026beyond,S23_yin2026contextualized,S02_figueiredo2025scaffolded}, support contextual interaction, whereas we examine whether the current encoded state and disclosure policy permit a proposed transition. Modeling affect and engagement \cite{S24_yannakakis2011experiencedriven,S25_melhart2020engagement,S26_wang2026engagement} targets different outcomes: style preferences, affect estimates, and personality labels may guide generation but do not authorize state mutations.

Existing benchmarks and automated play-testing cover state, embodiment, content, process, social interaction, and long-horizon play \cite{S11_cote2019textworld,S12_wang2022scienceworld,S13_fan2022minedojo,S14_perezliebana2019gvgai,S15_liu2024agentbench,S16_ma2024agentboard,S17_zhou2024sotopia,S18_paglieri2025balrog,S19_lepelletier2022playtesting}, while NCP-Bench evaluates narrative commitment and fact consistency \cite{S58_ma2026ncpbench}. Structural, schema, affordance, and neuro-symbolic rule methods motivate our constraint design \cite{S09_geng2023grammar,S10_beurerkellner2023prompting,S36_scholak2021picard,S50_kambhampati2024llmmodulo,S51_ichter2023saycan,S52_wang2025agentspec,S53_li2025settingdc,S54_zhou2025story2game,S42_buongiorno2024pangea,alshiekh2018shield,rebedea2023nemo}, alongside counterexample-guided repair \cite{S45_madaan2023selfrefine,S48_shinn2023reflexion,S46_chen2023selfdebug,S47_solarlezama2006sketching} and reproducibility/accounting standards \cite{S31_agarwal2021precipice,S33_pineau2021reproducibility,S56_feng2025agentrr,S59_sun2025bdc,gebru2021datasheets,akhtar2024croissant}. Unlike prior systems that conflate schema validity with action legality or treat factual grounding as sufficient for dialogue safety, \system{} isolates schema parsing, epistemic knowledge, disclosure permission, and canonical state commit into separately auditable gates.
